%% =============================================================================
%% SEC 1 - INTRODUCTION. ~620 words.
%% SOURCE: chapters/paper1_sec1_introduction.tex, compressed ~6:1.
%% v2 fix: stress-multiplier claim restated as in the paper (3-7x across the
%% nine strategy-benchmark combinations).
%% =============================================================================

\citet{duarte2007risk} ask whether fixed-income arbitrage earns genuine alpha
or merely picks up ``nickels in front of a steamroller.'' We take their
question to European markets and to three strategies that have each been
documented as a pricing anomaly but never examined as a systematic trading
strategy: the inflation-linked basis on Italian BTP~Italia bonds, the European
corporate CDS--bond basis, and the iTraxx index skew---the gap between the
traded index spread and the spread implied by its single-name constituents.
% SOURCE: sec1 opening paragraph, verbatim logic.

Three features make these trades unusually clean laboratory material. First,
each has a deterministic payoff at maturity: held to maturity, the profit is
locked in regardless of the interim path---unlike the model- and
assumption-dependent strategies of \citet{duarte2007risk}. Second, all three
are implementable: entry and exit follow mechanical basis thresholds,
transaction costs are charged at the trade level and widen with market stress,
and every reported return is net of those costs. Third, and central to this
paper, the three trades are run by the same type of leveraged institutional
arbitrageur but differ in who \emph{holds} the underlying instruments: the
CDS--bond and iTraxx legs sit on institutional balance sheets, whereas the
BTP~Italia is, by design of the Italian Treasury, held by unlevered retail
households who face no margin calls---the natural experiment on slow-moving
capital \citep{duffie2010presidential} that Section~\ref{sec:smc} exploits.
% SOURCE: sec1 contribution list items 1 and 3, merged.

Relative to the literature, the contribution is threefold. First, the
BTP~Italia basis has, to our knowledge, never been studied: work on linker
mispricing spans TIPS and index-linked gilts, but no academic paper covers the
BTP~Italia. It is also the harder puzzle---the TIPS--Treasury basis narrowed
as arbitrage capital flowed into the trade \citep{fleckenstein2014tips},
whereas the Italian one is still open thirteen years on, and it carries none
of the ballooning risk that complicates the TIPS trade, its two legs having
identical sovereign exposure. Second, while all three anomalies are documented
as pricing facts \citep{hull2004relationship, longstaff2005corporate,
junge2015liquidity}, none has been examined as an implementable trading
strategy with entry rules, exits, and stress-dependent transaction costs---the
step that turns a documented anomaly into a measurable, net-of-cost premium.
Third, where \citet{fleckenstein2014tips} document co-movement across three
U.S.\ mispricings that share an institutional holder base, our cross-section
varies the holder base at constant arbitrageur type---which is what separates
a common capital pool from a common risk factor, and turns slow-moving
capital from an assumption into a testable prediction.

Our findings, in brief. Net of costs, the strategies earn annualized Sharpe
ratios of 0.98 (BTP~Italia), 1.56 (CDS--bond), and 1.25 (iTraxx skew), with
positive skewness---the opposite of the steamroller caricature---and maximum
drawdowns of 3.0\%, 4.4\%, and 2.5\%, each recovered within months. The alpha---2.2\% to 4.8\% per year under the hardest control
layer---survives every layer of an escalating battery of risk
controls: the benchmark models of \citet{duarte2007risk},
\citet{fung2004hedge}, and \citet{brooks2020active}; eight principal
components of a 75-factor tradable universe; and best-subset selection that is
free to pick, for each strategy, whichever eight factors explain it best. The
selected models absorb up to 31\% of return variation and none of the alpha,
which also survives when every hedge is re-estimated out of sample from
past-only loadings. The alpha is state-dependent: the high-stress alpha is
three to seven times the calm-market alpha across the nine strategy--benchmark
combinations. And the slow-moving capital predictions hold with the sign
structure the holder-base experiment demands: the institutional pair co-moves
ever more tightly as stress builds (correlation rising from $-0.21$ to
$+0.32$), while the retail-held BTP~Italia moves \emph{against} the
institutional pair, reaching $-0.47$ in high stress; a widening of the
Libor--OIS spread pushes the two institutional mispricings wider and the
BTP~Italia basis tighter.
% SOURCE: sec1 results paragraphs; numbers from Tables I, XII, XIV, XIX.

The remainder of the paper follows this arc: Section~\ref{sec:strategies}
constructs the strategies and measures performance;
Section~\ref{sec:alpha} runs the alpha through five layers of controls;
Section~\ref{sec:smc} tests the slow-moving capital mechanism;
Section~\ref{sec:implications} draws the implications for investors.
Construction details, the full factor universe, and all robustness exhibits
are in the Appendix.
